实验显示该对象具有明显的慢性和空间分布特征：$T_2$ 比 $T_3$ 响应更快、幅值更大。采用 $K = 16.33$、$T = \SI{24}{\second}$ 的 PT1 模型可以用较少参数描述主要动态，并使代数控制器设计可行。仿真中的超调量近似满足 $5\,\%$ 要求，调节时间与所选 $\SI{100}{\second}$ 方案处于同一量级。
模型并不能覆盖所有真实效应。$T_1$ 在实验过程中明显变化，说明入口条件不恒定。另一方面，长时间 $\SI{3}{\volt}$ 数据、静态特性、$\SI{5}{\volt}$ 阶跃和独立闭环 dSpace 数据没有作为单独文件提供。本报告已完整说明这些任务应如何处理，并严格避免填入缺乏原始数据支持的测量数值。
